\bookchapter{第二十七章：爸妈也能轻松用——适老化手机完全指南}

\chapterintro{说句实话，写这一章的时候，我想到了我爸妈。我妈每次手机弹出一个“系统更新”，她都不敢点，怕把手机弄坏了。我爸的微信字体比蚂蚁还小，但他不知道怎么调大，就这么眯着眼睛看了两年。这不是他们“笨”，是手机默认把所有人都当成了二十多岁的年轻人。这一章，我用最直白的话，把长辈们手机使用中最头疼的问题一个个解开。}

\booksection{零、先把心态摆正：这不是“老人用不来”，是“手机没做好”}

在做任何设置之前，有一个观念要先说清楚：\textbf{现在的智能手机界面，从设计之初就没怎么考虑老年用户}。弹窗小得看不清、操作手势需要精准触碰、图标长得都差不多——这些不是用户的问题，是产品的问题。所以带长辈用手机，第一件事不是教，是帮他们\textbf{把环境调好}。环境调好了，大半的问题就消失了。

带长辈用手机有三条铁律：(1)\textbf{别催}——你用了十年的智能机，他们可能才用了三个月。(2)\textbf{用类比}——“Wi-Fi就像家里的水管，开了水龙头才有水”，这种比喻比专业术语好理解一万倍。(3)\textbf{写下来}——口头教一遍他们可能忘了，你拿张纸把关键步骤写下来贴在他手机壳背面，比什么都管用。

\booksection{一、第一步：打开“简易模式”——手机变清爽了}

现在的手机基本都内置了一个“简易模式”（有的叫“关怀模式”“老年模式”“长辈模式”），打开之后，整个手机都会大变样：

\begin{itemize}
  \item \textbf{字体和图标变得很大}——再也不用戴老花镜看手机了。
  \item \textbf{桌面图标大幅精简}——只保留电话、短信、相机、相册、微信等最常用的几个App，其他全收起来。
  \item \textbf{操作变简单了}——很多花里胡哨的功能（负一屏、智能助手、浮窗等）默认关闭，界面主打一个清爽。
  \item \textbf{声音变大了}——来电铃声、通知音量、通话声音都自动调到最大。
\end{itemize}

\booksubsection{各品牌简易模式开启方法（直接照着操作就行）：}

\begin{center}
\begin{tabularx}{\linewidth}{|X|X|X|}
\hline
\textbf{品牌} & \textbf{模式名称} & \textbf{打开方法} \tabularnewline
\hline
iPhone & 辅助使用（没有单独简易模式，但可逐项调大） & 设置 > 显示与亮度 > 文字大小，拉到最大。再进"辅助功能 > 显示与文字大小"开"粗体文本""按钮形状" \tabularnewline
华为/荣耀 & 简易模式 & 设置 > 系统和更新 > 简易模式 > 开启 \tabularnewline
小米/红米 & 极简模式 & 设置 > 更多设置 > 极简模式 > 开启 \tabularnewline
OPPO/一加 & 简易模式 & 设置 > 特色功能 > 简易模式 \tabularnewline
vivo/iQOO & 简易桌面 & 设置 > 桌面、锁屏与壁纸 > 简易桌面 \tabularnewline
三星 & 简易模式 & 设置 > 显示 > 简易模式 \tabularnewline
\hline
\end{tabularx}
\end{center}

如果你用的是其他品牌，在手机“设置”顶部的搜索框里直接搜“简易”或“长辈”两个字，一般都能找到。找不到的话，按下面的“手动调大”方案一条条设置，效果也一样。

\booksection{二、如果不开简易模式，手动调也可以——“大字+大图标”五步走}

有些长辈不想用简易模式，觉得界面太“老气”。那就手动调以下五项，效果一样好：

\begin{enumerate}
  \item \textbf{字体调到最大}：设置 > 显示 > 字体大小，拖到最右边。顺便把"显示大小"（有的叫"显示缩放"）也调到最大，这样所有界面元素都跟着变大。
  \item \textbf{图标变大}：在主屏幕双指捏合（像缩小照片那样）进入桌面设置 > 桌面布局（或图标大小），选格子最少的布局（比如3×5），图标自然就大了。
  \item \textbf{打开"高对比度文字"}：让文字边缘更锐利，在"辅助功能"里可以找到这个开关。
  \item \textbf{关闭不必要的通知}：设置 > 通知，把除了电话、短信、微信之外的App通知全关掉。这样手机不会一会儿响一下，不会让人觉得烦。
  \item \textbf{把常用联系人放到桌面}：长按桌面空白处 > 添加小工具 > 联系人，把老伴、儿女、医生、紧急联系人的快捷拨号放到桌面上。点一下就拨出去，不用翻通讯录。
\end{enumerate}

\begin{tcbpractice}{动手：帮长辈"装修"一下手机}
\begin{enumerate}
  \item 打开简易模式（按上面的品牌表格操作）。
  \item 把电量数字显示打开（设置 > 电池 > 电量百分比），让他们随时知道还剩多少电。
  \item 调出通知栏，把蓝牙、NFC、飞行模式这些长辈用不着的快捷开关移除或藏起来。
  \item 在桌面上放四个快捷拨号：老伴、儿女、家庭医生（如果有）、社区/物业。
\end{enumerate}

\end{tcbpractice}

\booksection{三、防诈骗——这章请一定带着长辈一起看}

老年人是电信诈骗的重灾区。不是他们警惕性低，而是骗子的话术经过了精心设计，专门针对老年人的心理弱点。以下是我整理的最常见骗局类型和应对方法，建议你直接念给长辈听，或者打印出来贴在电话旁边。

\booksubsection{六大常见骗局（看到了立刻挂电话）：}

\begin{itemize}
  \item \textbf{(1) "你孙子出事了，赶紧汇钱"}——骗子会说"你孙子/孙女在外面闯祸了，被警察抓了，需要保证金"。记住：\textbf{任何要求转账才能"放人"的电话都是诈骗}。公检法不会电话要钱。
  \item \textbf{(2) "你的医保卡被停用了"}——短信里说医保卡异常，让你点链接"激活"。记住：\textbf{医保局、社保局不会发短信让你点链接}。遇到这种情况直接打12333（社保官方热线）核实。
  \item \textbf{(3) "高收益理财，稳赚不赔"}——微信群、短信里推荐各种"养老理财""国家扶持项目"，承诺高额回报。记住：\textbf{承诺年化收益超过5\%还"保本"的，99.9\%是骗局}。
  \item \textbf{(4) "我是你领导/老朋友，换号码了"}——冒充熟人借钱。记住：\textbf{打电话核实比什么都管用}，说"不方便接电话"的八成是骗子。
  \item \textbf{(5) "免费领鸡蛋/保健品讲座"}——线下免费活动诱导老人听"健康讲座"，然后卖天价保健品。记住：\textbf{天上不掉馅饼，免费活动背后都是生意}。
  \item \textbf{(6) "你的手机中毒了，需要远程修复"}——骗子让你下载一个远程控制App（如向日葵、TeamViewer），然后远程操控你的手机转走钱。记住：\textbf{绝对不要让任何人远程操控你的手机！}手机厂商、通信运营商、银行都不会给你打电话让你装远程控制软件。
\end{itemize}

\booksubsection{三道防线，立好之后能挡掉90\%的诈骗：}

\begin{enumerate}
  \item \textbf{第一道：手机自带拦截}——安卓手机打开"骚扰拦截"（在拨号App或手机管家App里），iPhone打开"静音未知来电"（设置 > 电话 > 静音未知来电）。
  \item \textbf{第二道：运营商拦截}——发送短信"KTFSR"到10086，免费开通移动的高频骚扰电话拦截；电信联通也有类似服务。打运营商客服电话让人工帮你开通就行。
  \item \textbf{第三道：国家反诈中心App}——在应用商店下载"国家反诈中心"App，注册并开启来电预警。这个App和公安的数据库联动，骗子号码一拨过来就能识别。
\end{enumerate}

\begin{tcbwarn}{给长辈的"防骗口诀"（打印出来贴电话旁边）}
\begin{itemize}
  \item \textbf{陌生电话要转账 挂掉！}
  \item \textbf{短信链接要你点 删掉！}
  \item \textbf{说高收益保本的 拉黑！}
  \item \textbf{让你装远程控制软件 直接挂！}
  \item \textbf{真警察不会电话办案更不会要钱！}
  \item \textbf{拿不准的事，先给儿女打个电话问问。}
\end{itemize}

\end{tcbwarn}

\booksection{四、远程协助——你在外地也能帮爸妈修手机}

有时候长辈遇到的问题，你在电话里描述半天他们也找不到对应的按钮。这时候“远程协助”功能就派上用场了——你可以直接在自己的手机上看到他们手机的屏幕，帮他们操作。各品牌的方案不一样：

\begin{itemize}
  \item \textbf{华为/荣耀}：自带"畅连"App，里面有"屏幕共享"功能。你和爸妈都装畅连，发起视频通话后点击"共享屏幕"，你就可以在他们屏幕上画圈标注，告诉你妈妈"点这里"。
  \item \textbf{小米/红米}：小米视频通话App里有"远程协助"，支持屏幕共享+语音通话同步。
  \item \textbf{OPPO/vivo}：内置"家庭空间"或"远程守护"，在设置里搜索"远程"就能找到。
  \item \textbf{通用方案（所有品牌都可用）}：下载"向日葵远程控制"App（免费版够用），你们各装一个。注册同一个账号或加好友，让爸妈把手机上显示的验证码念给你，你就能远程看到和操作他们的屏幕了。
  \item \textbf{iPhone用户}：FaceTime通话中支持"共享屏幕"。在FaceTime通话时，点击屏幕共享按钮，对方就能看到你的屏幕。但iOS安全限制无法远程操作，只能"指挥式教学"。
\end{itemize}

\booksection{五、微信的“适老改造”——跟长辈说这几句就够了}

微信是长辈们用得最多的App，但界面太密太杂。以下是几个立竿见影的调整：

\begin{itemize}
  \item \textbf{调大微信自己的字体}：微信 > 我 > 设置 > 通用 > 字体大小，拖到最大。注意，这个只影响微信内部的字体，不影响手机系统字体。
  \item \textbf{打开"关怀模式"}：微信 > 我 > 设置 > 关怀模式 > 开启。打开后字体和按钮都会变大，色彩对比度加强，还能一键"听文字消息"（点一下文字消息就能朗读出来）。
  \item \textbf{帮他们备注好重要联系人}：把儿女、亲戚的名字在通讯录里改成容易认的——比如"大儿子-小明""二女儿-小红"，让他们一打开微信就知道谁是谁。
  \item \textbf{教他们用"语音输入"而不是打字}：聊天框旁边的键盘上，按住空格键（安卓）或点麦克风图标（iPhone），直接说话就能变成文字发送。识别准确率已经非常高了，对不会拼音打字的老人来说是"质变级"的便利。
  \item \textbf{关闭"发现页"里不用的功能}：微信 > 我 > 设置 > 通用 > 发现页管理，把"朋友圈""视频号""直播""购物""游戏"里不用的全关掉，减少干扰。
\end{itemize}

\booksection{六、手机丢了怎么办？帮长辈做个“应急预案”}

预则立，不预则废。趁手机还没丢，先把这几件事做好：

\begin{enumerate}
  \item \textbf{确认"查找我的手机"已打开}：iPhone在设置 > 你的名字 > 查找 > 查找我的iPhone；安卓在设置搜"查找设备"。务必确认这个开关是亮着的。
  \item \textbf{把重要的密码和账号记在一张纸上}：Apple ID / 华为账号 / 小米账号的密码，锁在家里抽屉里一份，给信得过的儿女一份。丢了手机之后要登录这些账号来远程锁定和定位，如果忘了密码就麻烦了。
  \item \textbf{教长辈一句话}：\textbf{"手机丢了，第一时间给家人打电话。"}家人可以立刻登录iCloud或安卓的查找设备网页，远程标记为"丢失模式"，这样捡到的人也打不开、用不了。
\end{enumerate}

\booksection{七、无线耳机和助听器兼容——被忽视的好功能}

很多长辈有听力问题，但不知道手机可以跟助听器协同工作。iPhone和部分安卓手机支持MFi（Made for iPhone）助听器，可以直接把电话声音传到助听器里，比举着手机听筒清楚太多了。

\begin{itemize}
  \item \textbf{连接方法}：iPhone设置 > 辅助功能 > 助听设备，打开蓝牙，按照助听器的配对说明操作。
  \item \textbf{如果不用助听器}：在"辅助功能 > 音频/视觉"里打开"电话降噪"和"单声道音频"，能改善通话清晰度。
  \item \textbf{安卓}：在"辅助功能 > 听力增强"里有类似设置，各品牌位置略有差异，直接搜"助听"两个字。
\end{itemize}

\booksection{八、健康相关功能——手机不只是通讯工具}

\begin{itemize}
  \item \textbf{SOS紧急求助}：连续按电源键五次（大多数手机），手机会自动拨打紧急电话并给预设的紧急联系人发送定位。帮长辈提前设置好——iPhone："设置 > 紧急SOS"，安卓："设置 > 安全 > 紧急求助"。教长辈用之前先演习一遍，免得到时候紧张按不准。
  \item \textbf{跌倒检测}：Apple Watch有这个功能，检测到你摔倒且一段时间不动，会自动拨打急救电话并通知家人。有预算的话给长辈配一个，关键时刻能救命。
  \item \textbf{用药提醒}：iPhone的"健康"App内置用药提醒功能，设好药名、时间和剂量，到点就弹出提醒。安卓可以用"定时闹钟"或下载"吃药提醒"类App。
\end{itemize}

\begin{tcbsummary}{本章小结}
\begin{itemize}
  \item 先开简易模式或手动调大字体/图标，环境调好大半问题就没了
  \item 带长辈用手机三铁律：别催、用类比、写下来
  \item 六大常见骗局要反复提醒，三道防线立好能挡掉90\%诈骗
  \item 远程协助让你人在外地也能帮爸妈修手机
  \item 微信开"关怀模式"+"语音输入"是长辈最实用的两个功能
  \item SOS紧急求助和跌倒检测，关键时刻救命的配置值得花五分钟设好
\end{itemize}

\end{tcbsummary}
